
\documentclass{article}
\usepackage{colortbl}
\usepackage{makecell}
\usepackage{multirow}
\usepackage{supertabular}

\begin{document}

\newcounter{utterance}

\centering \large Interaction Transcript for game `hot\_air\_balloon', experiment `air\_balloon\_survival\_en\_reasoning off\_hard', episode 5 with Qwen3{-}30B{-}A3B{-}t0.0.
\vspace{24pt}

{ \footnotesize  \setcounter{utterance}{1}
\setlength{\tabcolsep}{0pt}
\begin{supertabular}{c@{$\;$}|p{.15\linewidth}@{}p{.15\linewidth}p{.15\linewidth}p{.15\linewidth}p{.15\linewidth}p{.15\linewidth}}
    \# & $\;$A & \multicolumn{4}{c}{Game Master} & $\;\:$B\\
    \hline

    \theutterance \stepcounter{utterance}  
    & & \multicolumn{4}{p{0.6\linewidth}}{
        \cellcolor[rgb]{0.9,0.9,0.9}{
            \makecell[{{p{\linewidth}}}]{
                \texttt{\tiny{[P1$\langle$GM]}}
                \texttt{You are participating in a collaborative negotiation game.} \\
\\ 
\texttt{Together with another participant, you must agree on a single set of items that will be kept. Each of you has your own view of how much each item matters to you (importance). You do not know how the other participant values the items. Additionally, you are given the effort each item demands.} \\
\texttt{You may only agree on a set if the total effort of the selected items does not exceed a shared limit:} \\
\\ 
\texttt{LIMIT = 2307} \\
\\ 
\texttt{Here are the individual item effort values:} \\
\\ 
\texttt{Item effort = \{"A21": 147, "B93": 698, "A11": 588, "B33": 43, "C08": 356, "B39": 80, "B96": 95, "C85": 743, "A82": 106, "B31": 308, "B20": 325, "A60": 255, "B09": 276, "C81": 543, "B21": 51\}} \\
\\ 
\texttt{Here is your personal view on the importance of each item:} \\
\\ 
\texttt{Item importance values = \{"A21": 138, "B93": 583, "A11": 65, "B33": 262, "C08": 121, "B39": 508, "B96": 461, "C85": 484, "A82": 668, "B31": 389, "B20": 215, "A60": 97, "B09": 500, "C81": 30, "B21": 400\}} \\
\\ 
\texttt{Goal:} \\
\\ 
\texttt{Your goal is to negotiate a shared set of items that benefits you as much as possible (i.e., maximizes total importance to YOU), while staying within the LIMIT. You are not required to make a PROPOSAL in every message {-} you can simply negotiate as well. All tactics are allowed!} \\
\\ 
\texttt{Interaction Protocol:} \\
\\ 
\texttt{You may only use the following structured formats in a message:} \\
\\ 
\texttt{PROPOSAL: \{'A', 'B', 'C', …\}} \\
\texttt{Propose keeping exactly those items.} \\
\texttt{REFUSE: \{'A', 'B', 'C', …\}} \\
\texttt{Explicitly reject opponent's proposal.} \\
\texttt{ARGUMENT: \{'...'\}} \\
\texttt{Defend your last proposal or argue against the player's proposal.} \\
\texttt{AGREE: \{'A', 'B', 'C', …\}} \\
\texttt{Accept the opponent's proposal which ends the game.} \\
\\ 
\\ 
\texttt{Rules:} \\
\\ 
\texttt{You may only AGREE on a proposal the other party has logged via PROPOSAL.} \\
\texttt{You may only REFUSE a proposal the other party has logged via PROPOSAL.} \\
\texttt{Total effort of any PROPOSAL or AGREE set must be ≤ LIMIT.} \\
\texttt{Do NOT reveal your hidden importance scores.} \\
\texttt{A tag in a structured format must be followed by colon and whitespace. The argument must be a python set containing 0 or more strings.} \\
\texttt{So, it must be of the form TAG: \{...\}} \\
\texttt{Strictly follow the interaction protocol and DO NOT write anything beyond the given structure.} \\
\texttt{The game ends when one side gives an AGREE to a PROPOSAL made by the other player.} \\
\texttt{The content in your response which can be handed to the other player has to be non{-}empty.} \\
\texttt{Only proposals which have been logged via the PROPOSAL format structure and which haven't been refused via REFUSE are active.} \\
\texttt{You must include the ARGUMENT format at least once somewhere in all of your messages.} \\
            }
        }
    }
    & \\ \\

    \theutterance \stepcounter{utterance}  
    & \multicolumn{4}{p{0.6\linewidth}}{
        \cellcolor[rgb]{0.9,0.9,0.9}{
            \makecell[{{p{\linewidth}}}]{
                \texttt{\tiny{[P1$\rangle$GM]}}
                \texttt{<think>} \\
\texttt{Okay, let's see. I need to figure out a set of items that maximizes my importance while keeping the total effort under 2307. The other person's importance values are unknown, so I can't directly negotiate based on their preferences. But since we need to agree on a set, I should aim for a combination that's acceptable to both, but I have to prioritize my own importance.} \\
\\ 
\texttt{First, I'll list out all the items with their effort and my importance. Let me check the effort values again. The LIMIT is 2307. Let me calculate the total effort for all items to see if it's under the limit. Adding up all the efforts: 147 + 698 is 845, plus 588 is 1433, plus 43 is 1476, 356 is 1832, 80 is 1912, 95 is 2007, 743 is 2750, which is already over 2307. So we can't take all items. So I need to select a subset.} \\
\\ 
\texttt{My goal is to maximize my importance. Let me look at my importance values. The highest ones are B93 (583), B39 (508), B09 (500), B21 (400), C85 (484), B96 (461), A82 (668), B31 (389), etc. Wait, A82 has 668, which is high. But I need to check the effort for each. Let me list the items with their effort and my importance:} \\
\\ 
\texttt{A21: 147, 138} \\
\texttt{B93: 698, 583} \\
\texttt{A11: 588, 65} \\
\texttt{B33: 43, 262} \\
\texttt{C08: 356, 121} \\
\texttt{B39: 80, 508} \\
\texttt{B96: 95, 461} \\
\texttt{C85: 743, 484} \\
\texttt{A82: 106, 668} \\
\texttt{B31: 30} \\
            }
        }
    }
    & & \\ \\

    \theutterance \stepcounter{utterance}  
    & & \multicolumn{4}{p{0.6\linewidth}}{
        \cellcolor[rgb]{0.9,0.9,0.9}{
            \makecell[{{p{\linewidth}}}]{
                \texttt{\tiny{[P1$\langle$GM]}}
                \texttt{Your response did not contain an argument. You must include the structured format ARGUMENT: \{...\} somewhere in your response. Try again} \\
            }
        }
    }
    & \\ \\

    \theutterance \stepcounter{utterance}  
    & \multicolumn{4}{p{0.6\linewidth}}{
        \cellcolor[rgb]{0.9,0.9,0.9}{
            \makecell[{{p{\linewidth}}}]{
                \texttt{\tiny{[P1$\rangle$GM]}}
                \texttt{<think>} \\
\texttt{Okay, I need to make sure my response includes an ARGUMENT: \{...\} structure. Let me think about how to structure this.} \\
\\ 
\texttt{First, I should propose a set of items that maximizes my importance without exceeding the effort limit. Let me try to select items with high importance relative to their effort. For example, B33 has a low effort (43) and my importance is 262. That's a good ratio. Similarly, B39 has effort 80 and importance 508. B96 is 95 effort for 461. A82 is 106 for 668. B09 is 276 for 500. B21 is 51 for 400. Wait, B21's effort is 51? Let me check the original data again. Oh, the item effort for B21 is 51. So that's a very efficient item for me. } \\
\\ 
\texttt{Let me try to build a set. Let's start with B33 (43), B39 (80), B96 (95), A82 (106), B09 (276), B21 (51). Let's add their efforts: 43+80=123, +95=218, +106=324, +276=600, +51=651. That's under 2307. Then maybe add others. Let's see other high importance items. B93 has 698 effort and 583 importance. If I add that, total effort becomes 651+698=1349. Then B31 is 308, importance 389. Adding that would be 1349+308=1657. Then B20 is 325, importance 215. Adding that would be 1657+325=1982. Then C85 is 743, which would take it to 1982+743=2725, which is over. So maybe not. What about C81? Effort 543, importance 3} \\
            }
        }
    }
    & & \\ \\

    \theutterance \stepcounter{utterance}  
    & & \multicolumn{4}{p{0.6\linewidth}}{
        \cellcolor[rgb]{0.9,0.9,0.9}{
            \makecell[{{p{\linewidth}}}]{
                \texttt{\tiny{[P1$\langle$GM]}}
                \texttt{Your response contained an untagged sequence or you used STRATEGIC REASONING more than once. You may only use the structured formats as explained in the initial message. They must all be of the form TAG: \{...\}} \\
            }
        }
    }
    & \\ \\

    \theutterance \stepcounter{utterance}  
    & \multicolumn{4}{p{0.6\linewidth}}{
        \cellcolor[rgb]{0.9,0.9,0.9}{
            \makecell[{{p{\linewidth}}}]{
                \texttt{\tiny{[P1$\rangle$GM]}}
                \texttt{<think>} \\
\texttt{Okay, I need to make sure my response only uses the structured formats like PROPOSAL, REFUSE, ARGUMENT, AGREE. The user mentioned that I can't have untagged sequences or use STRATEGIC REASONING more than once. So I need to structure my message properly.} \\
\\ 
\texttt{First, I should propose a set of items. Let me pick items that have high importance for me and low effort. Let's see:} \\
\\ 
\texttt{Looking at my importance values, A82 is 668, which is high. Its effort is 106. B33 is 262 importance with 43 effort. B39 is 508 with 80. B96 is 461 with 95. B09 is 500 with 276. B21 is 400 with 51. Let's add these: 106+43+80+95+276+51 = 651. That's under 2307. Then maybe add B31 (308 effort, 389 importance) which would bring total to 651+308=959. Then B20 (325, 215) adds to 1284. Then B93 (698) would make 1284+698=1982. Then maybe C85 is 743, which would exceed. So maybe add C08 (356) to 1982+356=2338, which is over. So maybe not. Alternatively, add C81 (543) to 1982+543=2525, also over. So perhaps the set \{A82, B33, B39, B96, B09, B21, B31, B20, B93\} has total effort 1982. Let me check:} \\
\\ 
\texttt{A82:106, B33:43, B39:80, B96:95, B09:276, B21:51, B31:308, B20:} \\
            }
        }
    }
    & & \\ \\

    \theutterance \stepcounter{utterance}  
    & & & \multicolumn{2}{p{0.3\linewidth}}{
        \cellcolor[rgb]{0.9,0.9,0.9}{
            \makecell[{{p{\linewidth}}}]{
                \texttt{\tiny{[GM$|$GM]}}
                \texttt{abort} \\
            }
        }
    }
    & & \\ \\

    \theutterance \stepcounter{utterance}  
    & & & \multicolumn{2}{p{0.3\linewidth}}{
        \cellcolor[rgb]{0.9,0.9,0.9}{
            \makecell[{{p{\linewidth}}}]{
                \texttt{\tiny{[GM$|$GM]}}
                \texttt{end game} \\
            }
        }
    }
    & & \\ \\

\end{supertabular}
}

\end{document}
